\documentclass[10pt,a4paper]{amsart}
\usepackage[margin=19mm]{geometry}
\usepackage{amsmath,amssymb,mathtools}
\usepackage[scr=rsfs,cal=euler]{mathalfa}
\usepackage{xfrac}
\usepackage[unicode,hidelinks,bookmarks=false]{hyperref}
\newcommand{\ie}{i.e.\ }
\newcommand{\cf}{cf.\ }
\newcommand{\resp}{respectively\ }
\newcommand{\Subsection}{\subsection}
\providecommand{\nobreakdash}{\nobreak\hbox{-}}
\setlength{\emergencystretch}{2em}
\multlinegap0pt
\usepackage{xcolor}
\usepackage{amssymb}
\newtheorem{lemma}{Lemma}
\newcommand{\UM}[2]{\mathsf{U}_{#2,#1}} % uniform matroid
\DeclareMathOperator{\rk}{rk}
\title{\vspace{-10mm}The Tutte polynomial determines the linear coefficient of the Ehrhart polynomial}
\author{Erik Panzer}
\date{Source: arXiv:2608.28585v1 (CC BY 4.0)}
\begin{document}
\maketitle

\noindent\emph{Source and licence.} \vspace{-10mm}The Tutte polynomial determines the linear coefficient of the Ehrhart polynomial. Version arXiv:2608.28585v1 (CC BY 4.0), distributed by arXiv under the Creative Commons Attribution 4.0 International licence, http://creativecommons.org/licenses/by/4.0/, as recorded in the arXiv licence metadata for this exact version and shown on the abstract page https://arxiv.org/abs/2608.28585v1. Full licence notice: \texttt{licenses/arxiv-2608.28585-CC-BY-4.0.txt}.\par\smallskip

\noindent\textit{Editorial scope.} Counted unit: the article's lemma lem:upper-noloops (source0.tex lines 328--337). The article's complete original proof of it is delivered with it (source0.tex lines 338--350, one \texttt{proof} environment of 13 lines). No mathematics is written, repaired or re-derived: the statement and the proof are the article's own lines, in the article's own notation, and no retained line is substituted or re-flowed.  Retained uncounted as the source-local context the proof needs: lines 276--278; lines 319--324.  Nothing was omitted from any retained span, so the package carries no omission record.  No retained line contains a citation, so the wrapper carries no bibliography and no bibliography entry.  No retained line contains a figure environment, an image-inclusion command or drawing code, so no image is delivered and the package declares no asset dependency.  Editorial formatting only, in addition: an AMS \texttt{amsart} preamble that restates the article's own declarations and macros used by the retained lines, namely the environments \texttt{lemma} on separate counters, together with the standard packages those lines need.  The retained environments print in this document as Lemma 1 in order of appearance, whereas the article prints the same items with the numbering of its own full document.  The complete native source of the article as distributed in the arXiv e-print for this exact version is bundled unchanged as \texttt{sources/208798/source0.tex}.
\begin{lemma}\label{lem:ak-upper}
    For any matroid of size $n$ and rank $r$, we have $a_k(M)\leq r\binom{n-1}{k}$ for all $k$. Furthermore, an equality $a_k(M)=r\binom{n-1}{r}$ at any $0<k<n$ implies that $M\cong \UM{r}{r}\oplus\UM{\ell}{0}$.
\end{lemma}

\begin{lemma}\label{lem:contract-delete}
    If $e\in M$ is neither a loop nor a coloop of $M$, then for all $k>0$ we have that
    \begin{equation*}
        a_k(M)=a_k(M\setminus e)+a_{k-1}(M/e)\,.
    \end{equation*}
\end{lemma}

\begin{lemma}\label{lem:upper-noloops}
    Let $M$ be a loopless matroid of size $n$ and rank $r\geq1$. Then for all $k\geq0$,
    \begin{equation*}
        a_k(M)\leq a_k(\UM{r-1}{r-1}\oplus \UM{\ell+1}{1})
        %= (r-1)\binom{n-1}{k} + \binom{r-1}{k}
        = r\binom{n-1}{k} -\binom{n-1}{k} + \binom{r-1}{k}
        \,.
    \end{equation*}
    Furthermore, equality at any $2\leq k\leq n-1$ implies that $M\cong\UM{r-1}{r-1}\oplus\UM{\ell+1}{1}$.
\end{lemma}

\begin{proof}
    At $r=1$ the only loopless matroid is the bond, so $M\cong \UM{n}{1}\cong \UM{0}{0}\oplus \UM{\ell+1}{1}$. In this case, $\rk(A)=r=1$ for every non-empty subset $A\subseteq M$, hence $a_k(M)=0$ for all $k>0$ and $a_0(M)=1$. Thus $a_k(M)=\binom{0}{k}=(r-1)\binom{n-1}{k}+\binom{r-1}{k}$ for all $k$.
    
    At $r=n$, again there is only a single matroid $M\cong\UM{n}{n}\cong\UM{r-1}{r-1}\oplus\UM{1}{\ell+1}$ and the claim reduces to lemma~\ref{lem:ak-upper}, because then $(r-1)\binom{n-1}{k}+\binom{r-1}{k}=r\binom{n-1}{k}$.

    For $2\leq r\leq n-1$ there are several isomorphism classes of matroids and we proceed by induction on $r$. Since $\ell>0$ in this case, $M$ must have an element $e$ which is not a coloop. We can thus apply lemma~\ref{lem:contract-delete} to $e$. Since deletion cannot create loops, the induction hypothesis applies to $a_k(M\setminus e)$ and together with lemma~\ref{lem:ak-upper} we obtain
    \begin{align*}
        a_k(M) &\leq a_k(\UM{r-1}{r-1}\oplus\UM{\ell}{1}) + a_{k-1}(\UM{r-1}{r-1}\oplus\UM{\ell}{0})
        \\ &= (r-1)\binom{n-2}{k} + \binom{r-1}{k} + (r-1)\binom{n-2}{k-1}
        \\ &= (r-1)\binom{n-1}{k} + \binom{r-1}{k}.
    \end{align*}
    This proves the claimed bound. To get equality at some $2\leq k\leq n-1$ we must in particular have $a_{k-1}(M/e)=a_{k-1}(\UM{r-1}{r-1}\oplus\UM{\ell}{0})$, so lemma~\ref{lem:contract-delete} shows that $M/e\cong \UM{r-1}{r-1}\oplus\UM{\ell}{0}$. Contraction cannot create coloops, so the $r-1$ coloops of $M/e$ must already be present in $M$, that is, $M\cong\UM{r-1}{r-1}\oplus N$ for some matroid $N$ containing $e$. Since $M$ and thus $N$ have no loops, and $N/e\cong\UM{\ell}{0}$ has only loops, $N\cong\UM{\ell+1}{1}$ must be the bond. This proves $M\cong \UM{r-1}{r-1}\oplus\UM{\ell+1}{1}$.
\end{proof}

\end{document}
